\section{Experimental design}
\label{sec:experimental-design}

The experiments are designed to support a scoped claim about sparse nonnegative weighted digraphs, not a universal claim across all feedback arc set regimes. The sparse graph-benchmark collection provides the primary claim-bearing setting. Exact and MIP validation provide certified reference points on subsets where certification is feasible. EXP10 evaluates repeated-run behavior against an external executable whose initialization varies across runs. LOLIB serves as a transfer test that marks the dense complete-ordering boundary. Tables~\ref{tab:baseline-provenance} and~\ref{tab:comparison-accounting} summarize the comparison surface.

\begin{table}[t]
\centering
\tiny
\caption{Baseline provenance and comparison scope.}
\label{tab:baseline-provenance}
\setlength{\tabcolsep}{3pt}
\begin{tabularx}{\linewidth}{@{}p{1.35cm}p{1.05cm}p{1.25cm}p{1.35cm}p{1.35cm}cp{2.0cm}p{1.25cm}@{}}
\toprule
Method & Role & Origin & External / in-repo & Problem model & Weighted & Implementation behavior & Evaluation subset \\
\midrule
LR-TA & internal seed & \cite{DF03,VahidiKoutis2024arxiv} & in-repo & sparse FAS / ordering & yes & deterministic for fixed input and tie-breaking & sparse, exact, holdout \\
WMSF-style seed & internal seed & adapted from \cite{CC25,CCP24} & in-repo adaptation & sparse FAS & yes & deterministic for fixed input & sparse, exact, LOLIB \\
IPSNS & proposed method & this study & in-repo & sparse FAS refinement & yes & stochastic but reproducible for fixed configuration and seed & sparse, exact, holdout, LOLIB \\
Borda / net score & calibration & \cite{CFR10} & in-repo & ordering score rule & yes & deterministic & sparse, LOLIB \\
Weighted Eades adaptation & adaptation & adapted from \cite{ELS93,EL95} & in-repo adaptation & sparse ordering heuristic & yes & deterministic & sparse, LOLIB \\
Random multistart & calibration & standard multistart & in-repo & ordering heuristic & yes & stochastic & sparse, LOLIB \\
igraph Eades & external baseline & \cite{python_igraph_feedback_arc_set} & external maintained implementation & graph-library FAS heuristic & yes & deterministic as evaluated & sparse, LOLIB \\
DRMacIver/FAS & external baseline & \cite{drmaciver_feedback_arc_set} & external executable & matrix ordering & yes & evaluated executable shows run-to-run variation under internal initialization & sparse, EXP10, LOLIB, EXP9 \\
Exact DP & exact validator & repository implementation & in-repo & sparse FAS, \(n\le 20\) & yes & deterministic & exact validation \\
HiGHS MIP & exact validator & \cite{HuangfuHall2018} via \texttt{scipy.optimize.milp} & external maintained implementation & sparse-digraph MIP & yes & deterministic under fixed settings & medium MIP \\
\bottomrule
\end{tabularx}
\end{table}

\begin{table}[t]
\centering
\tiny
\caption{Comparison accounting for the main experiments.}
\label{tab:comparison-accounting}
\setlength{\tabcolsep}{3pt}
\begin{tabularx}{\linewidth}{@{}p{1.45cm}p{1.8cm}p{1.15cm}p{1.6cm}p{1.0cm}p{0.95cm}p{1.25cm}p{1.35cm}@{}}
\toprule
Experiment & Purpose & Instances attempted & Common / completed subset & Runs per method & Timeout & Primary metric & Statistical unit \\
\midrule
Primary sparse (EXP4) & main sparse comparison & 97 std.\ nonneg. & 97 internal; 93 paired with DR & 1 & wrapper; 4 DR incomplete & BW & instance \\
Exact DP (EXP3) & small-instance certification & 57 & 57 certified optima & 1 & \(n\le 20\) & gap to exact optimum & instance \\
MIP validation (EXP8) & medium-instance certification & 15 & 7 proven-optimal, 8 time-limited & 1 & 120\,s & gap when certified & instance \\
Holdout check & default selection & 43 total & 25 holdout instances & 5 & none & final BW by configuration & instance/config aggregate \\
Repeated-run (EXP10) & typical repeated-run behavior & 93 & 93 common sparse instances & 20 & EXP4 wrappers & median BW & paired instance difference \\
Dense transfer (EXP5) & complete-ordering boundary & 50 LOLIB & 50 & 1 & standard wrappers & BW & instance \\
\bottomrule
\end{tabularx}
\end{table}


\subsection{Benchmarks and scope}

The primary benchmark family comes from the public \texttt{graph-benchmarks} collection \cite{graph_benchmarks_repo}. After path-level deduplication, it yields 105 weighted sparse DIMACS digraph instances. The main benchmark tables restrict attention to the 97 standard instances whose aggregated arc weights are nonnegative. Eight instances are excluded from the standard set because they contain negative-weight arcs:
\texttt{gerez}, \texttt{howard-max}, \texttt{k3\_3}, \texttt{ku}, \texttt{peterson}, \texttt{peterson1}, \texttt{peterson2}, and \texttt{stg0}. This is a methodological scope restriction rather than a favorable filter: the local-ratio construction and the interpretation of incumbent protection used here are stated for nonnegative weights.

Three benchmark subsets play supporting roles. First, the exact-validation subset consists of the 57 standard nonnegative instances with \(n\le 20\), on which bitmask dynamic programming can certify the optimum. Second, a medium-scale subset of 15 instances supports the time-capped MIP validation. Third, a parameter study splits 43 sparse instances into 18 tuning instances and 25 holdout instances; its purpose is only to check whether the chosen IPSNS defaults are defensible engineering choices.

The dense transfer benchmark is a 50-instance subset of LOLIB 2010 \cite{MRD12,lolib_library}, comprising 25 SGB, 10 IO, and 15 RandA1 instances. These are complete weighted ordering instances rather than sparse general digraphs. They are included to delimit regime dependence: sparse-benchmark success should not be interpreted automatically as dense-ordering dominance.

\subsection{Methods compared}

IPSNS is the primary method under study. It starts from the better of two internal seeds and applies incumbent-protected SCC-local refinement. The two seeds are LR-TA and a WMSF-style adaptation of the weighted minimal/stable line \cite{CC25,CCP24}. They are retained in the tables because they clarify what the refinement adds and because the IPSNS non-worsening statement is relative to those seeds.

External and calibration baselines are chosen to cover accessible comparison classes rather than to exhaust the ranking literature. DRMacIver/FAS \cite{drmaciver_feedback_arc_set} is an external matrix-based ordering heuristic and the main external comparator in the sparse study. The evaluated executable exhibits run-to-run variation under internal initialization, which motivates EXP10. The python-igraph interface \cite{python_igraph_feedback_arc_set} supplies a graph-library Eades-style baseline. Weighted Eades, Borda/net-score, and random multistart are included as transparent local calibrators. Exact DP is used only on the 57-instance validation subset, and HiGHS MIP is used only for the medium-instance certification study.

\subsection{Evaluation metrics and accounting rules}

Backward weight (BW) is the primary metric throughout. Every method is evaluated under the same BW objective on the original aggregated arc weights, regardless of whether the method internally operates on an ordering, a removed arc set, or a matrix representation.

The sparse and dense comparison tables also report \emph{Best}, meaning the number of instances on which a method attains the minimum observed BW among the evaluated methods. Ties are credited to all tied methods. Relative excess over IPSNS is the mean percentage reduction achieved by IPSNS under the baseline-normalized convention \((X-\mathrm{IPSNS})/X\times 100\), where \(X\) is the comparator BW on the completed paired instances. Exact-validation tables instead report optimal-match counts and a mean gap normalized by total arc weight, matching the generator; the lone nonzero instance is discussed explicitly in the results text. Runtime is reported only as contextual wall-clock information and is not normalized across different external implementations.

Common-subset accounting is enforced wherever completion differs across methods. In the primary sparse comparison, DRMacIver/FAS completes 93 of the 97 standard instances; the remaining four incompletions remain visible rather than being hidden inside aggregate tables. The EXP10 repeated-run study therefore uses the 93-instance common subset completed by both IPSNS and DRMacIver/FAS. The MIP study reports certified gaps only on the 7 instances proved optimal within the time limit and treats the remaining 8 time-limited cases as incomplete reference evidence rather than heuristic failures.

\subsection{Statistical treatment}

Paired comparisons use the instance as the statistical unit. For EXP4 and EXP10, the paired outcome is the per-instance difference in BW between two methods on the relevant common subset. Archived summaries report two-sided Wilcoxon signed-rank and sign tests; the implemented Wilcoxon call uses SciPy's zero-difference removal convention (\texttt{zero\_method="wilcox"}), and the sign test excludes ties from the win/loss count. For EXP10 this leaves 38 nonzero paired median differences and 55 zero differences. Bootstrap confidence intervals are paired, use seed 42, and use 10{,}000 resamples. Relative excess in EXP10 is computed from per-instance medians using the same convention as EXP4:
\begin{equation}
    (\mathrm{DR}-\mathrm{IPSNS})/\mathrm{DR}\times 100.
    \label{eq:exp10-rel-excess}
\end{equation}
The archived summaries document the two-sided tests and bootstrap configuration directly; they do not materialize the Wilcoxon implementation mode, so the manuscript does not over-specify it.

\subsection{Repeated-run and holdout protocols}
\label{subsec:exp10-protocol}

EXP10 is a confirmatory repeated-run comparison on the 93-instance common sparse subset. Each instance receives 20 IPSNS runs with explicit integer seeds and 20 DRMacIver/FAS repetitions through the same wrapper and BW recomputation pipeline used in the main sparse comparison. The primary analysis compares per-instance median BW values with tie tolerance \(10^{-9}\). The protocol and analysis plan were fixed before examining aggregate outcomes. Reported summaries include wins, ties, losses, paired mean and median differences, mean relative excess, test results, effect size, and bootstrap intervals. Full per-instance summaries appear in Online Resource~1.

The holdout study is intentionally modest. It is not a claim-bearing benchmark comparison but a conservative check on engineering defaults. Across 18 tuning and 25 holdout sparse instances, six parameter configurations were each run with five seeds per instance, for 1290 archived runs in total. The holdout aggregate did not separate the tested defaults on median final BW, so the chosen IPSNS settings are interpreted as holdout-supported defaults rather than tuned optima.

\subsection{Implementation environment and reproducibility scope}

All reported numbers are taken from archived experiment outputs and regenerated manuscript assets. The runs used a Linux workstation with an Intel Core i7-12700K CPU, 62\,GiB RAM, Ubuntu~24.04, and Python~3.12. External tools were invoked through documented interfaces or wrappers, and failures or incompletions were retained explicitly in the summaries. Online Resource~1 supplies the deferred protocols, configurations, processed outputs, proofs, and regeneration scripts. Full reruns additionally require the documented external data, binaries, and computational resources.
